\section{Questions and Hypotheses}
\label{sec:questions}

The study has four research questions.

\begin{description}
  \item[RQ1: Correctness and control.] Do all four modes return the independent
  BFS cost, and do the three modes that ultimately use \HFull{} have identical
  successor-expansion traces?

  \item[RQ2: Mechanism.] Relative to \LazyFull{}, how many pivot evaluations,
  distance-table reads, and suffix calls does \Staged{} avoid, and how many
  additional OPEN operations does it introduce?

  \item[RQ3: Search time.] What is the paired staged/lazy search-time ratio
  over map-disjoint evaluation queries, with maps rather than timing
  repetitions treated as the top-level independent units?

  \item[RQ4: Regimes and amortization.] How do the mechanism and timing ratio
  vary across maze, random, room, and warehouse maps, and how does one-time
  landmark construction affect total cost at different query volumes?
\end{description}

The prospectively encoded analysis maps these questions to three hypotheses.
H1 is a hard gate: all 800 sealed queries must pass BFS equality and
full-landmark trace invariance before timing is analyzed.  H2 expects staged
evaluation to save exact pivot evaluations and distance reads relative to
ordinary lazy evaluation.  H3 evaluates the paired time ratio and the
association between saved work and time; it is explicitly descriptive.  In
particular, the bootstrap fraction of replicates with a ratio below one is
not treated as a frequentist $p$-value, and the Spearman coefficient does not
support a causal claim.
